\section{Elementary Flatness, Local Confluence, and Termination}
\label{sec:core-theorem}

Throughout this section $(X,\to)$ is an arbitrary abstract rewriting system. No finiteness is assumed, and every result has a Lean counterpart listed in \Cref{app:lean-formal-core}.

\subsection{The local equivalence}

\begin{lemma}[Reflexive joinability]
\label{lem:joinable-refl}
Every state is joinable with itself.
\end{lemma}

\begin{proof}
Take $d = x$ and use $x \to^{*} x$ twice.
\end{proof}

\begin{theorem}[Elementary flatness is local confluence]
\label{thm:elementary-flatness-iff-local-confluence}
An abstract rewriting system is elementarily flat if and only if it is locally confluent.
\end{theorem}

\begin{proof}
Suppose the system is locally confluent and let $(a,b,c)$ be an elementary peak. Then $a \to b$ and $a \to c$, so $b \joinable c$ by local confluence. Hence every elementary peak has a filler.

Conversely, suppose the system is elementarily flat and let $a \to b$ and $a \to c$. If $b = c$, then $b \joinable c$ by \Cref{lem:joinable-refl}. If $b \neq c$, then $(a,b,c)$ is an elementary peak, which has a filler by assumption, so again $b \joinable c$.
\end{proof}

\begin{remark}[What the theorem says]
\label{rem:core-theorem-scope}
Elementary flatness is local confluence restricted to peaks with distinct branches, so the theorem is close to a restatement; its only extra content is the trivial case $b = c$. We state it as a theorem because it is the exact form of the dictionary between the two vocabularies: ``every elementary branching has trivial holonomy'' and ``every one-step divergence can be rejoined'' are the same condition. The theorem says nothing about which peaks need to be checked in a structured formalism; that is the role of critical pairs (\Cref{sec:critical-pairs-curvature}). Nor is it, by itself, a local-to-global statement. Termination supplies the global step.
\end{remark}

\subsection{From local to global under termination}

The remaining results relate the local condition to confluence and to normal forms. The first is standard.

\begin{proposition}[Confluence and normalization give a unique normal form]
\label{prop:confluence-normalization-unique-normal-form}
If $(X,\to)$ is confluent and normalizing, then every $a \in X$ reaches exactly one normal form.
\end{proposition}

\begin{proof}
Normalization gives a normal form $n$ with $a \to^{*} n$. If also $a \to^{*} m$ with $m$ a normal form, confluence gives a $d$ with $n \to^{*} d$ and $m \to^{*} d$. A normal form has no descendants other than itself, so $d = n$ and $d = m$; hence $n = m$.
\end{proof}

Uniqueness alone, the ``at most one'' half of this argument, uses only confluence. The next lemma is the familiar fact that termination yields normal forms. We use the formulation of termination as well-foundedness: every nonempty set of states contains an element with no one-step successor inside the set. This permits induction along $\to$, in which a property may be assumed for all one-step successors of $a$ when proving it for $a$.

\begin{lemma}[Termination implies normalization]
\label{lem:terminating-normalizing}
A terminating abstract rewriting system is normalizing.
\end{lemma}

\begin{proof}
By well-founded induction along $\to$. If $a$ is a normal form, it reaches itself. Otherwise $a \to b$ for some $b$; by the induction hypothesis $b$ reaches a normal form $n$, and so $a \to b \to^{*} n$.
\end{proof}

\begin{theorem}[Newman's lemma]
\label{thm:newman}
A terminating, locally confluent abstract rewriting system is confluent.
\end{theorem}

\begin{proof}
We prove by well-founded induction along $\to$ that every $a$ satisfies: whenever $a \to^{*} b$ and $a \to^{*} c$, we have $b \joinable c$. If either sequence is empty, say $a = b$, then $c$ itself is a common descendant of $b$ and $c$. Otherwise both sequences begin with a step, $a \to b_1 \to^{*} b$ and $a \to c_1 \to^{*} c$. By local confluence there is a $d$ with $b_1 \to^{*} d$ and $c_1 \to^{*} d$. Since $a \to c_1$, the induction hypothesis applies at $c_1$; it joins $c$ and $d$ at some $v$. Since $a \to b_1$ and $b_1 \to^{*} d \to^{*} v$, the induction hypothesis applies at $b_1$; it joins $b$ and $v$ at some $w$. Then $b \to^{*} w$ and $c \to^{*} v \to^{*} w$.
\end{proof}

\begin{figure}[t]
\centering
\begin{tikzpicture}[x=1.0cm,y=1.0cm]
  \node[st] (a) at (0,0) {$a$};
  \node[st] (b1) at (-1.2,-1.0) {$b_1$};
  \node[st] (c1) at (1.2,-1.0) {$c_1$};
  \node[st] (d) at (0,-2.0) {$d$};
  \node[st] (b) at (-2.6,-2.2) {$b$};
  \node[st] (c) at (2.6,-2.2) {$c$};
  \node[st] (v) at (1.3,-3.2) {$v$};
  \node[st] (w) at (-0.4,-4.3) {$w$};
  \begin{scope}[on background layer]
    \fill[cellfill] (a.center)--(b1.center)--(d.center)--(c1.center)--cycle;
    \fill[blue!5] (c1.center)--(d.center)--(v.center)--(c.center)--cycle;
    \fill[blue!16] (b1.center)--(b.center)--(w.center)--(v.center)--(d.center)--cycle;
  \end{scope}
  \draw[rw] (a)--(b1); \draw[rw] (a)--(c1);
  \draw[rws] (b1)--(b); \draw[rws] (c1)--(c);
  \draw[rws] (b1)--(d); \draw[rws] (c1)--(d);
  \draw[rws] (d)--(v); \draw[rws] (c)--(v);
  \draw[rws] (b)--(w); \draw[rws] (v)--(w);
  \node[font=\scriptsize] at (0,-1.0) {LC};
  \node[font=\scriptsize] at (1.3,-2.2) {IH at $c_1$};
  \node[font=\scriptsize] at (-1.0,-2.9) {IH at $b_1$};
\end{tikzpicture}
\caption{The tiling in the proof of Newman's lemma. Local confluence (LC) fills the elementary peak at $a$; the induction hypothesis (IH), available because $b_1$ and $c_1$ are one step below $a$, fills the two larger regions. Termination is what makes the induction well founded.}
\label{fig:newman-tiling}
\end{figure}

\Cref{fig:newman-tiling} displays the proof as a tiling: one elementary $2$-cell and two cells supplied by induction. This is the precise sense in which, for terminating systems, global flatness is assembled from elementary flatness. Combining \Cref{thm:newman} with \Cref{thm:elementary-flatness-iff-local-confluence}, and noting that confluence trivially implies local confluence, we obtain the global form of the dictionary.

\begin{corollary}[Confluence is elementary flatness under termination]
\label{cor:terminating-confluence-iff-flat}
A terminating abstract rewriting system is confluent if and only if it is elementarily flat.
\end{corollary}

The relation between confluence and normal forms needs less than termination.

\begin{theorem}[Confluence and unique normal forms]
\label{thm:normalizing-confluence-iff-unf}
Let $(X,\to)$ be normalizing. Then $(X,\to)$ is confluent if and only if every state has unique reachable normal forms.
\end{theorem}

\begin{proof}
If the system is confluent, two normal forms reachable from the same state are joinable and hence equal, as in the proof of \Cref{prop:confluence-normalization-unique-normal-form}. Conversely, suppose every state has unique reachable normal forms, and let $a \to^{*} b$ and $a \to^{*} c$. By normalization, $b \to^{*} n$ and $c \to^{*} m$ for normal forms $n$ and $m$. Both are reachable from $a$, so $n = m$, and this common normal form joins $b$ and $c$.
\end{proof}

\begin{corollary}[Existence and uniqueness of normal forms]
\label{cor:terminating-flat-unique-nf}
In a terminating, elementarily flat abstract rewriting system, every state reaches exactly one normal form.
\end{corollary}

\begin{proof}
By \Cref{cor:terminating-confluence-iff-flat} the system is confluent, by \Cref{lem:terminating-normalizing} it is normalizing, and \Cref{prop:confluence-normalization-unique-normal-form} applies.
\end{proof}

For terminating systems, then, four conditions coincide: elementary flatness, local confluence, confluence, and uniqueness of reachable normal forms. Each hypothesis used along the way is needed. \Cref{subsec:boundary-search} exhibits smallest finite witnesses: a four-state system that is locally confluent but not confluent, for which termination fails (\Cref{fig:newman-needs-termination}); a one-state system that is confluent but not normalizing; and a three-state system in which every state has unique reachable normal forms but which is not confluent, for which normalization fails.

\begin{remark}[Mechanization]
\label{rem:lean-formal-core-crossref}
All results in this section are proved in Lean~4 for an arbitrary relation on an arbitrary type, with no finiteness assumption and with termination stated as well-foundedness of the reduction relation in the direction of rewriting. The proofs use no axioms beyond Lean's standard ones (propositional extensionality, quotient soundness, and choice), and several use none. \Cref{app:lean-formal-core} gives the correspondence.
\end{remark}
